\section{工程落地清单：面向 Agent 系统的设计原则}

\subsection{从课程内容抽象出的十条实践原则}
为了把讲者的系统观点转化为团队可执行方法，本节整理出面向 Agent 平台建设的十条工程原则：
\begin{enumerate}
    \item 把任务链路拆成可观测模块，避免黑箱联动；
    \item 先定义恢复策略，再上线高并发流量；
    \item 在系统设计早期就引入拓扑感知调度；
    \item 数据面设计与算力规划同步推进；
    \item 检索与重排要分层，先控成本再控准确；
    \item 企业场景必须做事实来源追踪；
    \item 把 SLA 与成本指标纳入同一优化目标；
    \item 把 ``模型升级'' 与 ``系统升级'' 解耦；
    \item 为边缘部署预留轻量模型路径；
    \item 持续做演练：故障恢复、流量峰值、链路退化。
\end{enumerate}

\begin{importantbox}{行动建议}
若团队资源有限，应优先补 ``可观测性 + 恢复能力 + 数据管线'' 三件事。这三者是把实验系统推向生产系统的最短路径。
\end{importantbox}

\subsection{评估框架}
\begin{longtable}{p{2.4cm}p{3.7cm}p{3.6cm}p{3.2cm}}
\toprule
\textbf{维度} & \textbf{评估问题} & \textbf{观测指标} & \textbf{常见改进手段} \\
\midrule
可靠性 & 故障后多久恢复？是否二次崩溃？ & MTTR、错误峰值、重试风暴次数 & 流控、预热、分级降级 \\
成本效率 & 单任务成本是否可预测？ & Token 成本、GPU 利用率、单位任务毛利 & 分阶段推理、缓存、模型路由 \\
数据质量 & 检索事实是否可追溯？ & 引用覆盖率、事实冲突率 & 数据版本化、检索评测 \\
可扩展性 & 高并发下是否退化平滑？ & P95/P99 延迟、队列积压长度 & 并发控制、弹性策略 \\
部署灵活性 & 能否同时支持云端与边缘？ & 端侧延迟、功耗、模型体积 & 蒸馏、量化、异构调度 \\
\bottomrule
\end{longtable}

\begin{knowledgebox}{为什么评估框架必须 ``系统化''}
Agent 质量不是单指标问题。只看准确率会忽略稳定性与成本，只看成本会牺牲可靠性。系统化评估能减少局部最优导致的整体失效。
\end{knowledgebox}

\subsection{本章小结}
课程中的方法论可以直接转化为工程清单：先建立系统能力，再追求模型能力红利，才能实现可持续迭代。

